\chapter{Граница голономного нуля и объёмной локализации}

Предыдущая глава получила одну инвариантную нулевую ветвь из
трёхциклической голономии, но её собственная функция оказалась постоянной
вдоль границы. Теперь проверяется последний прямой мост к физическому
дефекту: определяют ли уже построенные проекторная суперкривизна,
корневая связь и оператор переноса поперечный профиль без нового масштаба.

\section{Обратная сверка с прежней дефектной ветвью}

В проекте уже присутствуют два объекта, которые нельзя молча
отождествлять. Семейное поле
\begin{equation}
 H=\frac{2}{\sqrt3}\left(P-\frac14I\right)
\end{equation}
является нейтральным самосопряжённым эндоморфизмом. Его суперкривизна
\begin{equation}
 Q(H)=H^2-\frac{H}{\sqrt3}-\frac14I
 \label{eq:v5-localization-qh}
\end{equation}
выбирает четыре тетраэдрические оси, а гессиан
$\Tr Q(H)^2$ на диагональном бесследовом касательном пространстве имеет
три собственных значения $8/3$.

Майорановское поле спаривания $\Phi$, напротив, должно быть комплексным
заряженным сечением. Для вихря от него требуются ненулевой модуль вдали от
ядра, фазовая намотка и нуль в сердцевине. Ранние гейты вводили для него
независимые данные
\begin{equation}
 Z_\Phi,\qquad m_\Phi^2,\qquad\lambda_\Phi,\qquad g_\Phi.
 \label{eq:v5-localization-free-data}
\end{equation}

Эта разница не является вопросом обозначений. Поле $H$ кодирует семейную
ось и не несёт фазы заряда два; поле $\Phi$ несёт именно эту фазу.
Гейт калибровочно-семейного запирания также рассматривал их как различные
поля и доказал, что запирание не фиксирует знак гессиана $\Phi$.

\begin{nogocriterion}[Запрет переноса гессиана между полями]
Число $8/3$, выведенное для вариаций семейного эндоморфизма $H$, нельзя
объявить коэффициентом радиального потенциала $\Phi$, пока не построено
ковариантное отображение между нейтральным семейным сектором и заряженным
полем спаривания. Такое объявление было бы новым связующим коэффициентом.
\end{nogocriterion}

\section{Масштабный аудит радиального профиля}

Даже если условно предположить существование единичного вихря, его
минимальная радиальная энергия имеет вид
\begin{equation}
 E_\perp=2\pi\int_0^\infty r\,dr\left[
 Z_\Phi\left((\partial_r f)^2+\frac{f^2}{r^2}\right)
 +\frac{\lambda_\Phi}{2}(f^2-v^2)^2
 \right].
 \label{eq:v5-radial-gl-energy}
\end{equation}
Топология задаёт только граничные условия
\begin{equation}
 f(0)=0,\qquad f(\infty)=v,
\end{equation}
но не значения коэффициентов в \eqref{eq:v5-radial-gl-energy}.

Положим
\begin{equation}
 f(r)=vF(\rho),\qquad
 \rho=\frac r\xi,\qquad
 \xi=\sqrt{\frac{Z_\Phi}{\lambda_\Phi v^2}}.
 \label{eq:v5-coherence-length}
\end{equation}
Тогда уравнение для $F$ становится безразмерным, а энергия получает общий
множитель $2\pi Z_\Phi v^2$. Это может определить универсальную форму
$F(\rho)$ после выбора действия, но физическая ширина $\xi$ остаётся
свободной. Машинный аудит предъявляет два действия с одинаковыми
граничными условиями и одинаковым безразмерным профилем, но с
$\xi=1$ и $\xi=2$.

Следовательно, голономия фиксирует намотку, но не метрический размер
сердцевины. Оператор переноса вдоль дефекта также не устраняет эту
свободу: его дираковский предел фиксирует продольную кинематику, а не
норму поперечной производной $Z_\Phi$.

\section{Условная нулевая мода не исправляет масштаб}

Это видно уже в простейшем поперечном операторе Дирака с заданной вручную
стенкой
\begin{equation}
 m(x)=M\tanh\frac{x}{\xi}.
\end{equation}
Его нулевая мода пропорциональна
\begin{equation}
 \psi_0(x)\propto
 \cosh\!\left(\frac{x}{\xi}\right)^{-M\xi}.
 \label{eq:v5-domain-wall-zero}
\end{equation}
При $M\xi=1$ нормированная функция равна
\begin{equation}
 \psi_0(x)=\frac1{\sqrt{2\xi}}\operatorname{sech}\frac{x}{\xi},
\end{equation}
а её обратный коэффициент участия и среднеквадратичная ширина равны
\begin{equation}
 \operatorname{IPR}=\frac1{3\xi},\qquad
 \sqrt{\langle x^2\rangle}=\frac{\pi\xi}{\sqrt{12}}.
 \label{eq:v5-localization-width}
\end{equation}
Удвоение $\xi$ удваивает ширину и вдвое уменьшает обратный коэффициент
участия, хотя нормируемость и нечётная вещественная чётность ядра не
меняются.

Таким образом, топологическое утверждение защищает существование нуля
условно на заданном вихре, но не определяет его физический размер. Кроме
$\xi$ остаётся независимой комбинация $g_\Phi v$, задающая фермионную
длину затухания.

\begin{theorem}[Незамкнутость перехода от границы к ядру]
Трёхциклическая голономия и проекторная суперкривизна сохраняют одну
инвариантную граничную ветвь и семейную ось, но не выводят заряженное поле
спаривания, его ненулевой вакуум и поперечную кинетическую норму. Поэтому
они не определяют ни радиальный профиль, ни ширину локализованной нулевой
моды. Переход от граничного нуля к объёмному дефекту при текущем составе
полей не замыкается.
\end{theorem}

\section{Что именно сохранено}

Отрицательный результат не отменяет полученную структуру. Без параметров
сохраняются спектральные ветви $0,\pm1/3$, одномерная инвариантная линия и
условная нечётная майорановская чётность одного хирального канала. Кроме
того, старый вихревой результат остаётся корректной условной теоремой:
если заряженный конденсат и его вихрь уже построены, нулевая мода может
быть нормируемой.

Не получено другое: единый родительский функционал, который одновременно
порождает $\Phi$, фиксирует $Z_\Phi$, выбирает $v\ne0$ и связывает
$g_\Phi$ с тем же следом. До появления такого функционала нельзя называть
граничный нуль локализованным нейтрино или физической частицей.

\section{Критерий переоткрытия}

Мост допускается переоткрыть только после предъявления одного
предварительно фиксированного родительского построения, в котором:
\begin{enumerate}
 \item поле $\Phi$ возникает как правильно заряженное сечение, а не как
 переименование $H$;
 \item его ковариантная кинетическая норма и потенциал происходят из одного
 следа или одной меры;
 \item ненулевой вакуум и вихревой профиль следуют из вариационного
 уравнения;
 \item фермионная связь фиксируется тем же построением;
 \item длина локализации вычисляется до сравнения с наблюдаемыми.
\end{enumerate}

При отсутствии этих данных дальнейший подбор функций $m(r)$, ширин или
портальных коэффициентов запрещён. Следующий шаг должен не усложнять
профиль, а заморозить дефектно-транспортную ветвь и сформулировать точный
тип отсутствующего заряженного родительского поля.

\begin{protocol}
\begin{enumerate}
 \item граничная нулевая ветвь: \textbf{сохранена};
 \item нейтральное проекторное поле $H$: \textbf{выведено};
 \item заряженное поле спаривания $\Phi$: \textbf{не выведено};
 \item каноническое отождествление $H\leftrightarrow\Phi$:
 \textbf{отсутствует};
 \item безразмерная форма профиля: \textbf{условно возможна};
 \item физическая ширина $\xi$: \textbf{не фиксирована};
 \item нормируемая объёмная мода: \textbf{условна на заданном вихре};
 \item локализованный нейтринный дефект: \textbf{не получен};
 \item физическое замыкание: \textbf{не достигнуто}.
\end{enumerate}
\end{protocol}

Машинный протокол сохранён в
\path{s2t_v5_holonomy_zero_mode_localization_gate_results.json}.